\documentclass[11pt,a4paper]{article}
\usepackage[margin=2.4cm]{geometry}
\usepackage{amsmath,amssymb}
% Deliberately no hyperref / parskip / booktabs: they are not installed here
% and would trigger a package-installer prompt on every build. DOIs are
% therefore plain text.
\setlength{\parindent}{0pt}
\setlength{\parskip}{6pt}

\title{Grid Search: the four display statistics}
\author{pyHYSCORE / \texttt{grid\_search.py}}
\date{}

\newcommand{\SSE}{\mathrm{SSE}}
\newcommand{\RMSD}{\mathrm{RMSD}}

\begin{document}
\maketitle

The grid search scans two spin-Hamiltonian parameters and, at every grid
point, simulates each active dataset and compares it with experiment. The
\textbf{Display} pulldown chooses which of four statistics is mapped. This
note defines each one, explains what it does and does not tell you, and
gives references for the last three.

\tableofcontents

\section{Notation}

\subsection*{Indices and per-dataset quantities}

\begin{tabular}{@{}l p{12.6cm}@{}}
\hline
Symbol & Meaning \\
\hline
$\theta=(\theta_1,\theta_2)$ & the pair of scanned spin-Hamiltonian parameters; one grid point \\
$k$ & index over active (shown) datasets \\
$D$ & the data: all active datasets taken together \\
$e_k$ & experimental spectrum of dataset $k$, restricted to its region of interest, as a vector \\
$s_k(\theta)$ & simulation for dataset $k$ at $\theta$, interpolated onto the same points, as a vector \\
$N_k$ & number of spectrum points in dataset $k$'s region of interest \\
$f_{\min},f_{\max}$ & frequency limits of that region, the same on both axes \\
$c_k(\theta)$ & best-fit non-negative amplitude applied to $s_k$, Eq.\ \eqref{eq:amp} \\
$\SSE_k(\theta)$ & residual sum of squares for dataset $k$, Eq.\ \eqref{eq:sse} \\
$\mathrm{null}_k$ & null-model sum of squares $\lVert e_k\rVert^2$: the residual left by no simulation \\
$r_k$ & uncentred correlation between $e_k$ and $s_k$, Eq.\ \eqref{eq:sin} \\
$\sigma_k$ & unknown noise level of dataset $k$; a nuisance parameter, marginalized away \\
$N_k^{\text{eff}}$ & effective (independent) point count, $N_k/F$ \\
$a$ & an arbitrary rescaling factor applied to a dataset, used only in invariance arguments \\
\hline
\end{tabular}

\subsection*{Aggregate and displayed quantities}

\begin{tabular}{@{}l p{12.6cm}@{}}
\hline
Symbol & Meaning \\
\hline
$N$ & pooled point count $\sum_k N_k$ \\
$\RMSD(\theta)$ & display 1: pooled root-mean-square residual, Eq.\ \eqref{eq:rmsd} \\
$\RMSD_{\text{norm}}(\theta)$ & display 2: null-normalized score in $(0,1]$, Eq.\ \eqref{eq:norm} \\
$p(\theta\mid D)$ & posterior probability density over the scanned parameters, Eq.\ \eqref{eq:post} \\
$\Phi(\theta)$ & negative log posterior, the quantity actually minimized, Eq.\ \eqref{eq:nlp} \\
$i,j$ & indices of a cell on the interpolated (fine) display grid \\
$a,b$ & dummy cell indices inside sums over the grid \\
$p_{ij}$ & display 4: posterior normalized over the grid, $\sum_{ij}p_{ij}=1$, Eq.\ \eqref{eq:grid} \\
$L_{ij}$ & display 3: credible level of a cell, Eq.\ \eqref{eq:level} \\
$p^{(1)}_i,\;p^{(2)}_j$ & marginal distributions along each scanned axis \\
$F$ & oversampling: spectrum points per independent measurement, $F=4^{z}$ \\
$z$ & zero-fill setting from the main window \\
$\SSE_{\min},\;\RMSD_{\min}$ & smallest value of that quantity over the scanned grid \\
$\Delta$ & fractional rise in $\SSE$ above its minimum that bounds a confidence region, Eq.\ \eqref{eq:thresh}; $\Delta_F$ and $\Delta_{\chi^2}$ are its two prescriptions, Eqs.\ \eqref{eq:ftest} and \eqref{eq:chi2} \\
$\hat{\sigma}$ & noise level estimated from the fit itself, as opposed to known independently \\
$\nu$ & denominator degrees of freedom of the $F$ distribution, $\nu=N-p$ \\
$F_{p,\,N-p}(\text{conf})$ & quantile of the $F$ distribution with $p$ and $N-p$ degrees of freedom \\
$\chi^2_p(\text{conf})$ & quantile of the chi-square distribution with $p$ degrees of freedom \\
$\text{conf}$ & the confidence or credible fraction, here $0.68$ or $0.95$ \\
\hline
\end{tabular}

\textbf{One collision to be aware of.} $p$ carries two standard meanings in
this document, kept because both are conventional in their own literature:
as a function of arguments, $p(\cdot)$ is always a probability density; as a
bare symbol in Section \ref{sec:ftest} it is the number of jointly scanned parameters,
which here is always $2$. The subscripted $p_{ij}$, $p^{(1)}_i$, $p^{(2)}_j$
are probabilities.

\subsection*{Operators and abbreviations}

\begin{tabular}{@{}l p{12.6cm}@{}}
\hline
Symbol & Meaning \\
\hline
$\langle x,y\rangle$ & Euclidean inner product $\sum_n x_n y_n$ over the region of interest \\
$\lVert x\rVert$ & Euclidean norm, $\lVert x\rVert^2=\langle x,x\rangle$ \\
$\lceil x\rceil$ & ceiling: the smallest integer not less than $x$ \\
$\max(x,0)$ & clamp to zero: $x$ if positive, otherwise $0$ \\
\hline
RMSD & root-mean-square deviation \\
SSE & sum of squared residuals \\
roi & region of interest: the $[f_{\min},f_{\max}]^2$ window a fit is evaluated over \\
HPD & highest posterior density; a credible region containing the most probable values \\
FFT & fast Fourier transform \\
HYSCORE & hyperfine sublevel correlation spectroscopy \\
\hline
\end{tabular}

\section{What one grid point produces}

Write $\theta=(\theta_1,\theta_2)$ for the pair of scanned parameters. For
each active dataset $k$ let $e_k$ be the experimental spectrum restricted to
that dataset's region of interest $[f_{\min},f_{\max}]^2$, containing $N_k$
points, and let $s_k(\theta)$ be the simulation interpolated onto the same
grid. Both are treated as vectors in $\mathbb{R}^{N_k}$.

A simulated HYSCORE spectrum carries no meaningful absolute intensity, so
each dataset is given its own amplitude, fitted in closed form by least
squares and clamped to be non-negative:
\begin{equation}
c_k(\theta) \;=\; \max\!\left(\frac{\langle e_k,\,s_k(\theta)\rangle}
{\lVert s_k(\theta)\rVert^{2}},\; 0\right).
\label{eq:amp}
\end{equation}
The clamp matters later: a negative amplitude would mean an inverted
simulation fits best, which is unphysical here.

The residual sum of squares for that dataset is then
\begin{equation}
\SSE_k(\theta) \;=\; \bigl\lVert e_k - c_k(\theta)\,s_k(\theta)\bigr\rVert^{2}
\;=\;\sum_{\text{roi}} \bigl(e_k - c_k s_k\bigr)^2 .
\label{eq:sse}
\end{equation}
Everything below is a different way of combining the set
$\{\SSE_k, N_k\}$ into one number per grid point.

\section{Display 1 --- RMSD}

\begin{equation}
\RMSD(\theta) \;=\; \sqrt{\frac{\sum_k \SSE_k(\theta)}{\sum_k N_k}} .
\label{eq:rmsd}
\end{equation}

This is the pooled root-mean-square residual: all residuals from all
datasets are thrown into one pool and the root mean square is taken. It is
the original fitness measure and remains the default.

\textbf{What it means.} A distance in experimental intensity units.
Because those units are arbitrary (spectrometer gain, number of scans), the
absolute value carries no information --- only comparisons between grid
points within one scan are meaningful.

\textbf{The property to be aware of.} Equation \eqref{eq:rmsd} weights
datasets by their absolute intensity. If dataset $B$ is $250\times$ brighter
than dataset $A$, its residuals are $250\times$ larger and its contribution
to the sum is about $62\,500\times$ larger: $B$ effectively decides the fit
alone. This is a modelling choice, not a defect, but it is rarely the
intended one when several measurements of comparable quality are combined.

\section{Display 2 --- Normalized RMSD}
\label{sec:norm}

Define each dataset's \emph{null model}: the residual left by no simulation
at all,
\begin{equation}
\mathrm{null}_k \;=\; \lVert e_k \rVert^{2} \;=\; \sum_{\text{roi}} e_k^2 .
\end{equation}
The displayed statistic is the point-count-weighted mean of the per-dataset
unexplained fractions:
\begin{equation}
\RMSD_{\text{norm}}(\theta) \;=\;
\sqrt{\frac{\sum_k N_k \bigl(\SSE_k(\theta)/\mathrm{null}_k\bigr)}
{\sum_k N_k}} .
\label{eq:norm}
\end{equation}

\subsection{Why it is bounded}

Setting $c_k=0$ in \eqref{eq:sse} gives $\SSE_k=\lVert e_k\rVert^2 =
\mathrm{null}_k$ exactly. Since \eqref{eq:amp} clamps $c_k\ge 0$, the fitted
amplitude can never do \emph{worse} than $c_k=0$, so
\begin{equation}
0 \;<\; \frac{\SSE_k}{\mathrm{null}_k} \;\le\; 1
\qquad\Longrightarrow\qquad
0 \;<\; \RMSD_{\text{norm}} \;\le\; 1 .
\end{equation}
This is a hard bound, not an empirical observation. The interpretation is
absolute: $0$ is a perfect fit, $1$ means the simulation explains nothing.

\subsection{Geometric reading}

When $\langle e_k,s_k\rangle>0$, substituting the optimum \eqref{eq:amp}
into \eqref{eq:sse} gives
$\SSE_k = \lVert e_k\rVert^2 - \langle e_k,s_k\rangle^2/\lVert s_k\rVert^2$,
hence
\begin{equation}
\sqrt{\frac{\SSE_k}{\mathrm{null}_k}} \;=\; \sqrt{1-r_k^{2}},
\qquad
r_k \;=\; \frac{\langle e_k, s_k\rangle}
{\lVert e_k\rVert\,\lVert s_k\rVert},
\label{eq:sin}
\end{equation}
where $r_k$ is the uncentred correlation coefficient. The normalized value
is therefore the sine of the angle between experiment and simulation viewed
as vectors, and $r_k^2$ is the fraction of experimental power the simulation
accounts for. A displayed value of $0.3$ means $1-0.3^2 = 91\,\%$ of the
power is explained.

\subsection{Brightness invariance, and its cost}
\label{sec:brightness}

Rescaling a dataset, $e_k \to a\,e_k$, multiplies both $\SSE_k$ and
$\mathrm{null}_k$ by $a^2$, leaving the ratio unchanged. Every dataset
therefore contributes on equal terms regardless of signal strength.

The cost is that \eqref{eq:norm} is a genuinely different statistic from
\eqref{eq:rmsd}: it is a weighted sum of $\SSE_k$ with weights
$N_k/\mathrm{null}_k$ rather than uniform weights. With datasets of unequal
brightness the two surfaces have different shapes, and the best-fit point
can move between them. Neither is the display being wrong --- they answer
different questions about how much a faint dataset should count.

\subsection{Caveat}

The normalization is uncentred: no baseline is subtracted. Magnitude
spectra have a positive noise floor, which inflates $\mathrm{null}_k$ and
flatters the score. Two datasets with different baselines are therefore not
strictly comparable on this scale.

\section{Displays 3 and 4 --- posterior probability}

\subsection{From residuals to a probability}

Assume the residuals of dataset $k$ are Gaussian with an unknown noise
level $\sigma_k$, independently for each dataset. The likelihood is
\begin{equation}
p\bigl(D \mid \theta, \{\sigma_k\}\bigr) \;\propto\;
\prod_k \sigma_k^{-N_k}
\exp\!\left(-\frac{\SSE_k(\theta)}{2\sigma_k^{2}}\right).
\end{equation}
The noise levels are nuisance parameters. Assigning each the
scale-invariant Jeffreys prior $p(\sigma_k)\propto 1/\sigma_k$ and
integrating,
\begin{equation}
\int_0^{\infty}\!\sigma^{-N-1}
\exp\!\left(-\frac{\SSE}{2\sigma^{2}}\right)\mathrm{d}\sigma
\;\propto\; \SSE^{-N/2},
\end{equation}
which leaves a posterior over the scanned parameters alone:
\begin{equation}
\boxed{\;
p(\theta \mid D) \;\propto\; \prod_k \SSE_k(\theta)^{-N_k/2}\;}
\label{eq:post}
\end{equation}
(times a prior on $\theta$, taken flat over the scanned box). Equivalently,
the quantity minimized is
\begin{equation}
\Phi(\theta) \;=\; -\log p(\theta\mid D)
\;=\; \sum_k \frac{N_k}{2}\,\log \SSE_k(\theta) \;+\; \text{const}.
\label{eq:nlp}
\end{equation}

\subsection{Two consequences}

\textbf{It is the logarithm of each dataset's residual that gets averaged},
not the residual itself. That is what marginalizing an unknown per-dataset
noise level produces, and it is the principled answer to the weighting
question raised in Section \ref{sec:brightness} --- the weights are inferred rather than
chosen.

\textbf{It is automatically invariant to brightness and to the
normalization of Section \ref{sec:norm}.} Rescaling $e_k\to a\,e_k$ sends
$\log \SSE_k \to \log \SSE_k + 2\log a$, a constant shift of $\Phi$ that
cancels on normalization. Likewise
$\log(\SSE_k/\mathrm{null}_k) = \log \SSE_k - \log \mathrm{null}_k$ differs
by a constant. The probability map is therefore identical whether or not
the RMSD is normalized --- the question that makes displays 1 and 2 disagree
simply does not arise here.

\subsection{Effective number of points}
\label{sec:neff}

$N_k$ enters \eqref{eq:post} as an exponent, so it controls how sharp the
posterior is; the width of a credible region scales as $N^{-1/2}$.
Adjacent points of a processed spectrum are not independent: apodization
correlates neighbours, and zero-filling manufactures points carrying no new
information at all. The \textbf{Oversampling} control divides the point
count accordingly,
\begin{equation}
N_k^{\text{eff}} \;=\; \frac{N_k}{F},
\qquad
F \;=\; 4^{\,z} \;=\; \bigl(2^{\,z}\bigr)^2 ,
\end{equation}
where $z$ is the zero-fill setting: \texttt{update\_FFT} pads to
$2^{\lceil \log_2 N\rceil + z}$, so each axis gains $2^{z}$ points and the
two-dimensional count gains that squared. Leaving $F=1$ overstates the
information content and yields credible regions too narrow by a factor
$\sqrt{F}$.

\subsection{On the grid}

$\Phi$ is evaluated at the coarse scan points, interpolated (the negative
log posterior is smooth and roughly quadratic near its minimum, whereas the
probability itself spans many orders of magnitude and interpolates badly),
then exponentiated and normalized over the fine grid:
\begin{equation}
p_{ij} \;=\;
\frac{\exp\bigl(-[\Phi_{ij}-\min \Phi]\bigr)}
{\sum_{ab}\exp\bigl(-[\Phi_{ab}-\min \Phi]\bigr)},
\qquad \sum_{ij} p_{ij}=1 .
\label{eq:grid}
\end{equation}
Failed grid points get $p_{ij}=0$. Note that \eqref{eq:grid} makes the
scanned box the prior: the result is only meaningful if the posterior has
decayed to negligible before the edges, which the panel checks and warns
about.

\subsection{Display 3 --- credible level}

Each cell is labelled with the total probability mass of all cells at least
as probable as it:
\begin{equation}
L_{ij} \;=\!\!\sum_{\{ab\,:\; p_{ab}\,\ge\, p_{ij}\}}\!\! p_{ab}
\;\in\;(0,\,1].
\label{eq:level}
\end{equation}
This is the smallest highest-posterior-density region that contains the
cell. It runs from near $0$ at the most probable point to $1$ at the least,
so contouring at $0.68$ and $0.95$ draws exactly the $68\,\%$ and $95\,\%$
credible regions. Cells of equal density share one level, so a ring of
equally probable cells is never split across a contour.

This is the recommended view: it is readable everywhere, and its contours
are credible regions by construction rather than by a threshold rule.

\subsection{Display 4 --- posterior density}

The same object before the cumulative step: $p_{ij}$ from \eqref{eq:grid}
itself, the probability mass per grid cell. It is the primary quantity, and
$L$ is a monotone function of it, so both produce identical contours. Its
drawback is visual: a sharp posterior puts nearly every cell at effectively
zero and only a few pixels light up. It is most useful when the posterior
is broad.

\subsection{Uncertainty reported in these modes}

Instead of the F-test region used for displays 1 and 2, the interval comes
from the marginal distributions
\begin{equation}
p^{(1)}_i = \sum_j p_{ij},
\qquad
p^{(2)}_j = \sum_i p_{ij},
\end{equation}
reported as the highest-density interval containing $95\,\%$ of the marginal
mass. Marginalizing is the principled version of what the bounding box of a
joint region approximates, and it is generally tighter.

\section{Confidence intervals for displays 1 and 2}
\label{sec:ftest}

Displays 1 and 2 carry no noise model, so an interval cannot come from a
posterior. It comes instead from the shape of the residual surface around
its minimum, and the \textbf{CI method} dropdown offers two ways of turning
that shape into a region. This section derives both, shows how they are
related, and quantifies the difference --- which turns out to be negligible.

\subsection{What is being tested}

By construction $\SSE(\theta)\ge \SSE_{\min}$ everywhere. The question is
how much worse than the best fit a parameter pair may be while remaining
consistent with the data. Both methods answer it with a multiplicative
threshold,
\begin{equation}
\SSE(\theta) \;\le\; \SSE_{\min}\,\bigl[1+\Delta\bigr]
\qquad\Longleftrightarrow\qquad
\RMSD(\theta) \;\le\; \RMSD_{\min}\sqrt{1+\Delta},
\label{eq:thresh}
\end{equation}
and differ only in what they use for $\Delta$. The square root appears
because RMSD is the square root of a mean of $\SSE$; the code applies the
threshold in RMSD units directly.

\subsection{F-test: the noise level is unknown}

Since $\sigma$ is not known independently, it has to be estimated from the
fit itself, as $\hat{\sigma}^2 = \SSE_{\min}/(N-p)$. The excess sum of
squares is then compared against that estimate,
\begin{equation}
\frac{\bigl[\SSE(\theta)-\SSE_{\min}\bigr]/p}
{\SSE_{\min}/(N-p)} \;\sim\; F_{p,\,N-p},
\end{equation}
which is exact for a linear model and asymptotically valid otherwise.
Rearranging into the form of \eqref{eq:thresh},
\begin{equation}
\Delta_{F} \;=\; \frac{p}{N-p}\,F_{p,\,N-p}(\text{conf}).
\label{eq:ftest}
\end{equation}
The essential feature is the \emph{denominator} degrees of freedom $N-p$:
they encode the fact that $\hat{\sigma}$ is itself uncertain, having been
estimated from the same finite dataset. This is why the F-test is the
appropriate choice when the noise level is not known independently, which
is the situation here.

\subsection{Delta chi-square: the noise level is treated as known}

If $\sigma$ really were known, the excess sum of squares would be
$\chi^2$-distributed and the region would be
$\SSE(\theta)-\SSE_{\min}\le \sigma^2\chi^2_p(\text{conf})$. Substituting
the fit's own variance estimate $\hat{\sigma}^2=\SSE_{\min}/N$ in place of
$\sigma^2$ gives
\begin{equation}
\Delta_{\chi^2} \;=\; \frac{\chi^2_p(\text{conf})}{N}.
\label{eq:chi2}
\end{equation}
The substitution is the weak point: $\hat{\sigma}$ is treated as if it were
exact, so the uncertainty in the noise estimate is not propagated. The
result is therefore optimistic --- the region is always at least slightly
too small.

\subsection{How the two are related}

The connection is the limiting behaviour of the $F$ distribution. As the
denominator degrees of freedom grow,
\begin{equation}
p\,F_{p,\,\nu}(\text{conf}) \;\xrightarrow[\nu\to\infty]{}\; \chi^2_p(\text{conf}),
\end{equation}
which for $p=2$ and $\text{conf}=0.95$ reads
$2\times 2.9958 = 5.9915 = \chi^2_2(0.95)$. Substituting this limit into
\eqref{eq:ftest} gives $\Delta_{F}\approx \chi^2_p/(N-p)$, so the two
prescriptions differ in two small ways:
\begin{equation}
\Delta_{F}-\Delta_{\chi^2}
\;\approx\;
\chi^2_p\left(\frac{1}{N-p}-\frac{1}{N}\right)
\;=\;\frac{p\,\chi^2_p}{N(N-p)}
\;=\;\mathcal{O}\!\left(N^{-2}\right),
\end{equation}
plus the (also $\mathcal{O}(N^{-1})$) excess of $F_{p,\nu}$ over its own
limit at finite $\nu$. Both corrections have the same sign, so
\begin{equation}
\Delta_{F} \;>\; \Delta_{\chi^2} \quad\text{always},
\end{equation}
and the F-test is invariably the more conservative of the two. They agree
in the limit of infinite data.

\subsection{How much this actually matters}

For a residual surface that is locally quadratic, the half-width of the
interval scales as $\sqrt{\Delta}$, so the ratio
$\sqrt{\Delta_F/\Delta_{\chi^2}}$ is the factor by which choosing the
F-test widens the reported interval. At $\text{conf}=0.95$, $p=2$:

\begin{center}
\begin{tabular}{rrrrrr}
\hline
$N$ & $F_{p,N-p}$ & $1+\Delta_F$ & $1+\Delta_{\chi^2}$
& RMSD threshold (F / $\chi^2$) & width ratio \\
\hline
50      & 3.1907 & 1.132947 & 1.119829 & 1.064400 / 1.058220 & $1.053\times$ \\
100     & 3.0892 & 1.063045 & 1.059915 & 1.031041 / 1.029522 & $1.026\times$ \\
400     & 3.0184 & 1.015168 & 1.014979 & 1.007555 / 1.007461 & $1.006\times$ \\
1000    & 3.0047 & 1.006022 & 1.005991 & 1.003006 / 1.002991 & $1.003\times$ \\
2400    & 2.9995 & 1.002502 & 1.002496 & 1.001250 / 1.001247 & $1.001\times$ \\
10000   & 2.9966 & 1.000599 & 1.000599 & 1.000300 / 1.000300 & $1.000\times$ \\
\hline
\end{tabular}
\end{center}

At the point counts this program actually works with --- a region of
interest holds thousands of points --- the two methods differ by about one
part in a thousand in the reported interval. That is far below the
resolution of the scan grid, and well below the width of a single contour
line on the map. The distinction only becomes visible below $N\approx 100$,
which does not occur here.

\subsection{What to worry about instead}

The choice of test is swamped by two other effects.

\textbf{The effective point count.} $N$ appears directly in both
\eqref{eq:ftest} and \eqref{eq:chi2}, and the interval width scales roughly
as $N^{-1/2}$. Using the oversampling divisor $F$ of
Section \ref{sec:neff} at $\text{conf}=0.95$ and a raw $N=2400$:

\begin{center}
\begin{tabular}{rrr}
\hline
Oversampling $F$ & $N^{\text{eff}}$ & interval width, relative to $F=1$ \\
\hline
1  & 2400 & $1.00\times$ \\
4  & 600  & $2.01\times$ \\
16 & 150  & $4.06\times$ \\
64 & 38   & $8.57\times$ \\
\hline
\end{tabular}
\end{center}

A realistic zero-fill setting changes the interval by a factor of several,
against the one-part-in-a-thousand from the choice of test. Getting the
effective point count roughly right matters some three orders of magnitude
more than choosing between F and $\chi^2$.

\textbf{The shared assumptions.} Both methods assume independent, Gaussian,
zero-mean residuals, a correctly specified model, and a surface that is
approximately quadratic near the minimum. Real spectra violate all three to
some degree --- systematic error from orientation selection, blind spots and
baseline artifacts is not noise, and magnitude spectra have Rayleigh- rather
than Gaussian-distributed noise near the floor. Those violations affect the
interval far more than the difference between \eqref{eq:ftest} and
\eqref{eq:chi2}.

\textbf{Practical recommendation.} Leave the setting on F-test. It is the
correct one for an unknown noise level, it is the conservative one, and it
costs nothing. Then spend the attention on the oversampling factor and on
whether the residual maps show structure the model is failing to reproduce.

\subsection{Relation to the posterior displays}

The two frameworks are not in opposition. For a linear model with a
Jeffreys prior on $\sigma$, the marginal posterior over the parameters is a
multivariate $t$ distribution whose highest-density regions coincide exactly
with the F-test confidence regions. Displays 3 and 4 are therefore the same
statement as \eqref{eq:ftest} in the linear-Gaussian limit; where they
differ in practice is that the posterior uses the actual shape of
$\Phi(\theta)$ rather than assuming it is quadratic, and that it handles
several datasets with separate unknown noise levels correctly, which the
dataset-reweighted mean \eqref{eq:norm} only approximates.

\section{Summary}

\begin{center}
\begin{tabular}{llll}
\hline
Display & Range & Units & Weighting of datasets \\
\hline
RMSD                & $(0,\infty)$ & experimental intensity & by absolute intensity \\
Normalized RMSD     & $(0,1]$      & dimensionless          & equal (arithmetic) \\
Credible level      & $(0,1]$      & probability mass       & inferred (logarithmic) \\
Posterior density   & $(0,1]$      & mass per cell          & inferred (logarithmic) \\
\hline
\end{tabular}
\end{center}

\section{References}

The last three displays have close analogues in the structural-biology and
Bayesian-inference literature.

\subsection*{Normalized RMSD}

G. Cornilescu, J. L. Marquardt, M. Ottiger, A. Bax,
\emph{Validation of protein structure from anisotropic carbonyl chemical
shifts in a dilute liquid crystalline phase}, J. Am. Chem. Soc.
\textbf{120}, 6836--6837 (1998). \\
\texttt{doi:10.1021/ja9812610} \quad \texttt{https://doi.org/10.1021/ja9812610}

Introduces the quality factor
$Q=\mathrm{rms}(D_{\mathrm{obs}}-D_{\mathrm{calc}})/\mathrm{rms}(D_{\mathrm{obs}})$,
algebraically the same construction as \eqref{eq:norm}: residual divided by
the intrinsic scale of the data. Arrived at for the same reason --- the raw
number is in arbitrary units and means nothing on its own.

A. T. Br\"unger, \emph{Free R value: a novel statistical quantity for
assessing the accuracy of crystal structures}, Nature \textbf{355},
472--475 (1992). \\
\texttt{doi:10.1038/355472a0} \quad \texttt{https://doi.org/10.1038/355472a0}

The canonical treatment of normalized residual factors, and of why a fit
statistic evaluated on the data used to fit it is optimistic. The
cross-validation remedy transfers directly: hold out one $\tau$ value and
test the parameters found from the others.

\subsection*{Posterior probability (credible level and density)}

G. L. Bretthorst, \emph{Bayesian Spectrum Analysis and Parameter
Estimation}, Lecture Notes in Statistics \textbf{48}, Springer (1988). \\
\texttt{doi:10.1007/978-1-4684-9399-3} \quad
\texttt{https://doi.org/10.1007/978-1-4684-9399-3}

The closest reference of all. Derives \eqref{eq:post} for spectral data:
amplitude parameters that enter linearly are marginalized analytically, an
unknown noise level is integrated out under a Jeffreys prior, and what
remains is a posterior over the nonlinear parameters proportional to a
negative power of the residual sum of squares. Equations \eqref{eq:amp} and
\eqref{eq:post} are that scheme applied to a two-dimensional spectrum.

E. T. Jaynes (ed. G. L. Bretthorst), \emph{Probability Theory: The Logic of
Science}, Cambridge University Press (2003). \\
\texttt{doi:10.1017/CBO9780511790423} \quad
\texttt{https://doi.org/10.1017/CBO9780511790423}

The foundational treatment of nuisance-parameter marginalization and of
Jeffreys priors for scale parameters, which is what turns an unknown
$\sigma_k$ into the exponent $-N_k/2$.

W. Rieping, M. Habeck, M. Nilges, \emph{Inferential structure
determination}, Science \textbf{309}, 303--306 (2005). \\
\texttt{doi:10.1126/science.1110428} \quad
\texttt{https://doi.org/10.1126/science.1110428}

Applies exactly this idea to structure determination: rather than choosing
weights for heterogeneous data by hand, treat the error scale of each data
set as an unknown and marginalize it. This is the argument for why
\eqref{eq:nlp} is preferable to \eqref{eq:norm} when combining datasets.

M. Habeck, M. Nilges, W. Rieping, \emph{Replica-exchange Monte Carlo scheme
for Bayesian data analysis}, Phys. Rev. Lett. \textbf{94}, 018105 (2005). \\
\texttt{doi:10.1103/PhysRevLett.94.018105} \quad
\texttt{https://doi.org/10.1103/PhysRevLett.94.018105}

The sampling counterpart, for when a two-parameter grid is no longer
enough; also discusses inferring the weight of each data source.

\end{document}
